\chapter{Admissibility, Policies, and Sampling}
\label{ch:admissibility}

A generative system must produce histories that are admissible and, beyond that, that are good. This chapter studies the two conditions in turn. The first is existence: when is it possible to continue a story indefinitely without becoming inadmissible? The second is distribution: how should a probabilistic model be corrected so that it never leaves the admissible set, and in what sense is the correction optimal?

\section{Infinite admissible histories}

Let $\Adm\subseteq E^*$ be a prefix-closed set of admissible histories. A history $h\in\Adm$ is \emph{extendable} if $he\in\Adm$ for some event $e$. The set $\Adm$ is \emph{non-blocking} if every $h\in\Adm$ is extendable, so that a story is never forced to stop.

\begin{proposition}[Existence of infinite admissible chains]\label{prop:chains}
Suppose $\Adm$ is nonempty, prefix-closed, and non-blocking. Then there exists an infinite sequence of events $e_1,e_2,\dots$ all of whose finite prefixes lie in $\Adm$.
\end{proposition}

\begin{proof}
Choose $h_0=\varepsilon$, which lies in $\Adm$ because $\Adm$ is nonempty and prefix-closed. Given $h_k\in\Adm$, non-blockingness provides an event $e_{k+1}$ with $h_ke_{k+1}\in\Adm$; set $h_{k+1}=h_ke_{k+1}$. The sequence $(e_k)$ so obtained has every prefix $h_k$ in $\Adm$. The construction uses a choice at each step, that is, dependent choice, and no further assumption.
\end{proof}

The hypotheses are not vacuous. Non-blockingness can fail in a way that is easy to miss: a story with a linear resource that must be consumed, as in Chapter~\ref{ch:linear}, may reach a marking from which the acceptable terminal set cannot be reached. Such a history is a \emph{dead end}. The check for dead ends is a reachability problem for the underlying resource net, and a generative system that samples without lookahead will, with positive probability, enter them.

When the set of events is finite, a stronger conclusion holds under weaker hypotheses.

\begin{proposition}[K\"onig]\label{prop:konig}
Suppose $E$ is finite and $\Adm$ is prefix-closed and contains histories of arbitrarily large length. Then there exists an infinite sequence all of whose prefixes lie in $\Adm$.
\end{proposition}

\begin{proof}
Call $h\in\Adm$ \emph{good} if $\Adm$ contains arbitrarily long extensions of $h$. The empty history is good by hypothesis. If $h$ is good, then, since $E$ is finite and extensions of length greater than $\lvert h\rvert$ begin with $he$ for some $e\in E$, at least one $he\in\Adm$ is good, for otherwise the length of admissible extensions of $h$ would be bounded by the maximum of the bounds for the finitely many $he$. Choosing a good $e_{k+1}$ at each step yields the sequence.
\end{proof}

\section{Policies}

A \emph{policy} is a rule that selects the next event given the history. Formally it is a function $\pi:\mathcal{A}'\to E$ on a prefix-closed subset $\mathcal{A}'\subseteq\Adm$ with $h\cdot\pi(h)\in\mathcal{A}'$ for all $h\in\mathcal{A}'$, or more generally a function to probability distributions on $E$ supported on admissible events.

\begin{definition}[Admissible policy]
A policy $\pi$ is \emph{admissible} if the history generated by iterating $\pi$ from $\varepsilon$ lies in $\Adm$ at every finite stage, for every realization of the randomness.
\end{definition}

\begin{corollary}\label{cor:policy}
An admissible policy exists if and only if $\Adm$ contains a nonempty non-blocking prefix-closed subset. In particular, if $\Adm$ is non-blocking, the policy defined in the proof of \Cref{prop:chains} is admissible.
\end{corollary}

The practical content is that admissibility alone does not make a good generator. A policy that always chooses the first admissible event is admissible and typically dull. What distinguishes a good policy is a distribution over admissible continuations that also reflects style, surprise, and rhythm, and for that the language of probability is needed.

\section{Conditioning as optimal projection}

Let $P$ be a probability measure on a countable set $W$ of candidate continuations, for instance the output distribution of a trained model, and let $A\subseteq W$ be the admissible ones, with $P(A)>0$. The \emph{admissibility projection} of $P$ is the conditional measure
\[
P_A(w)=\frac{P(w)\,\mathbf{1}_A(w)}{P(A)}.
\]
It is realized in practice by rejection sampling: draw from $P$ and discard samples outside $A$.

\begin{theorem}[Optimality of conditioning]\label{thm:projection}
Among all probability measures $Q$ on $W$ with $\supp Q\subseteq A$, the measure $P_A$ is the unique minimizer of $\KL(Q\,\|\,P)$, and
\[
\KL(P_A\,\|\,P)=-\log P(A).
\]
\end{theorem}

\begin{proof}
For $Q$ supported on $A$ with $\KL(Q\|P)<\infty$, write
\[
\KL(Q\|P)=\sum_{w\in A}Q(w)\log\frac{Q(w)}{P(w)}
=\sum_{w\in A}Q(w)\log\frac{Q(w)}{P_A(w)}+\sum_{w\in A}Q(w)\log\frac{P_A(w)}{P(w)}.
\]
The first sum is $\KL(Q\|P_A)\ge0$, with equality if and only if $Q=P_A$, by Gibbs' inequality. In the second sum, $P_A(w)/P(w)=1/P(A)$ for $w\in A$, so it equals $-\log P(A)$ independently of $Q$. Therefore $\KL(Q\|P)=\KL(Q\|P_A)-\log P(A)\ge-\log P(A)$, with equality exactly when $Q=P_A$. Evaluating at $Q=P_A$ gives the stated value.
\end{proof}

The theorem identifies the cost of enforcing admissibility. If the model puts probability $P(A)$ on admissible continuations, enforcement through conditioning incurs a KL divergence of $-\log P(A)$ from the model. Rejection sampling needs on average $1/P(A)$ attempts per accepted sample. A model whose admissible mass is $2^{-20}$ will need on the order of a million attempts per accepted sample, and the cost of enforcing admissibility by rejection scales exactly as the reciprocal of the admissible mass. This is the quantitative argument for building admissibility into training or decoding, as opposed to filtering outputs afterwards.

\begin{remark}[Masking is not conditioning]\label{rem:masking}
Constrained decoding, which masks inadmissible next events at each step and renormalizes, differs in general from whole-sequence conditioning, even when $\Adm$ is prefix-closed and non-blocking. At a prefix $h$, conditioning on eventual admissibility selects the next event $e$ with probability proportional to $P(e\mid h)\,P(\Adm\mid he)$, where $P(\Adm\mid he)$ is the probability that a continuation of $he$ is admissible. Masking omits the factor $P(\Adm\mid he)$ and weights each locally admissible event by $P(e\mid h)$ alone. The two agree exactly when (assuming $P(e\mid h)>0$ for the locally admissible events) at every admissible prefix, this factor takes the same value for all locally admissible next events. A small example shows the difference: let $E=\{0,1\}$, let $P$ be uniform on words of length two, and let the admissible length-two words be $00$, $01$, $10$. This set is prefix-closed and non-blocking, since the prefix $1$ extends by $0$. Conditioning gives each admissible word probability $1/3$, so the first event is $0$ with probability $2/3$. Masking leaves both first events locally admissible and chooses each with probability $1/2$. The difference is the weight $P(\Adm\mid he)$ that lookahead supplies. When $\Adm$ can block, masking can also enter dead ends, which is a further reason lookahead is needed.
\end{remark}

\section{Where learned systems stand}

Current learned systems divide, broadly, into those that maintain an explicit latent state used for planning and control~\cite{ha2018worldmodels,hafner2020dreamer} and those that generate frames conditioned on recent frames and actions with no persistent state beyond a compressed context~\cite{valevski2024gamengen,bruce2024genie}. In the vocabulary of Chapter~\ref{ch:layers}, the first kind has a state layer and the second kind has an appearance layer with a short window of history. Neither, as described in the cited work, maintains the committed history as a separate authoritative object. The mathematical analysis suggests what such a layer would give: a set $\Adm$ against which every generated event can be checked and a fold from which state can be recomputed.

\section{Bridge: from valid continuation to cinematic decision}

Admissibility determines which events may follow, not how they are shown. In the forged-letter scene the reveal at $t=58$ is admissible under many realizations. Suppose the pause before the question can last any whole number of seconds from $0$ to $8$, the reaction shot any whole number of seconds from $0$ to $6$, and the music can begin at the door slam, two seconds after it, or not at all. There are $9\times7\times3=189$ realizations on this grid. All are admissible, they share a single history, and they differ in anticipation, emphasis, and the opportunities they give a viewer to settle. Choosing among them is a decision about timing, framing, performance, and sound that occurs after admissibility has been settled and is not made by it.

Where such a decision belongs in the layered world is a design question. A choice that changes only appearance, for instance the exact length of a reaction shot, can be discarded and regenerated without affecting what the world has committed to. A choice that alters what the audience has been told, or what a character can reasonably do next, changes the fold and must be recorded as an event in the history. The decision rule is thus the one stated in the authority principle: a realization choice enters the history exactly when it affects the semantic fold.

\section*{Exercises}

\begin{exercise}
Give an example of a prefix-closed admissible set that contains arbitrarily long histories but is not non-blocking, and show that \Cref{prop:konig} still produces an infinite chain when $E$ is finite. Does the example contradict \Cref{prop:chains}?
\end{exercise}

\begin{exercise}
Show that conditioning is idempotent: $(P_A)_A=P_A$. Show also that for $A'\subseteq A$ with $P(A')>0$ one has $(P_A)_{A'}=P_{A'}$.
\end{exercise}

\begin{exercise}
Suppose the model assigns admissible mass $P(A)=0.01$. How many samples are needed on average to obtain one admissible sample by rejection? Compute $\KL(P_A\|P)$ in bits.
\end{exercise}
